\begin{figure}[htbp]
	\centering
	\begin{tikzpicture}[font=\small]
		\def\W{4.6}\def\h{0.6}\def\a{0.85}\def\r{0.55}
		% concio
		\draw[trvconcio] (0,-\h) rectangle (\W,\h);
		\draw[trvfibra] (0,-\h-0.12) -- (\W,-\h-0.12);
		% faccia sinistra: azioni con verso opposto a quello positivo
		\draw[trvforza] (0,0) -- ++(-\a,0) node[left] {$N(z_1)$};
		\draw[trvforza] (0,0) -- ++(0,\a) node[above] {$T(z_1)$};
		\pic at (0,0) {trvmomgiu};
		\node[below left] at ($(0,0)+(235:0.42)$) {$M(z_1)$};
		% faccia destra: azioni nei versi positivi
		\draw[trvforza] (\W,0) -- ++(\a,0) node[right] {$N(z_2)$};
		\draw[trvforza] (\W,0) -- ++(0,-\a) node[right, pos=0.9] {$T(z_2)$};
		\pic at (\W,0) {trvmomsu};
		\node[above right] at ($(\W,0)+(55:0.42)$) {$M(z_2)$};
		% azioni concentrate in z_i
		\draw[trvforza] ({\W/2},0) -- ++(0.75,0) node[above=1pt, pos=0.7] {$Q_i$};
		\draw[trvforza] ({\W/2},0) -- ++(0,-0.5) node[right=1pt, pos=0.7] {$P_i$};
		\draw[trvcoppia] ({\W/2-0.75},-0.12) arc[start angle=0, end angle=180, radius=0.24];
		\node[below] at ({\W/2-0.99},-0.12) {$C_i$};
		\pic at ({\W/2},0) {trvpunto};
		% quote
		\draw[trvrichiamo] (0,-\h-0.2) -- (0,-1.6);
		\draw[trvrichiamo] ({\W/2},-\h-0.2) -- ({\W/2},-1.6);
		\draw[trvrichiamo] (\W,-\h-0.2) -- (\W,-1.6);
		\draw[trvquota] (0,-1.45) -- node[below=1pt] {$\tfrac{1}{2}\Delta z$} ({\W/2},-1.45);
		\draw[trvquota] ({\W/2},-1.45) -- node[below=1pt] {$\tfrac{1}{2}\Delta z$} (\W,-1.45);
		\node[below left=0pt] at (0,-1.45) {$z_1$};
		\node[below right=0pt] at (\W,-1.45) {$z_2$};
		\node[above] at ({\W/2},\h) {$z_i$};
	\end{tikzpicture}
	\caption{Concio di lunghezza $\Delta z$ centrato nella
		sezione di ascissa $z_i$, dove agiscono la forza di
		componenti $P_i$, $Q_i$ e la coppia $C_i$. Per
		$\Delta z \to 0$ le azioni interne sulle facce tendono
		ai limiti da sinistra e da destra in $z_i$, e
		l'equilibrio del concio fornisce le equazioni di
		salto~\eqref{eq:salto_NTM}.}
	\label{fig:concio_salto}
\end{figure}
